% automatisch erzeugt von tools/anhang_b.py -- nicht von Hand ändern

\subsubsection{Eingabe}
\offen{Satz: Verweis auf Tabelle~\ref{tab:h1_profil} (Organisationsprofil)}
\begin{small}
\begin{longtable}{@{}>{\raggedright\arraybackslash}p{5.0cm}>{\raggedright\arraybackslash}p{9.5cm}@{}}
\caption[Organisationsprofil für Fall~1]{Organisationsprofil für Fall~1. Quelle: eigene Darstellung, Daten aus \repodatei{h1/input/firmenprofil.json}.}\label{tab:h1_profil}\\
\toprule \textbf{Merkmal} & \textbf{Angabe} \\ \midrule \endfirsthead
\toprule \textbf{Merkmal} & \textbf{Angabe} \\ \midrule \endhead
\bottomrule \endfoot
Name & Hochschule München \\
Land & Deutschland \\
Bundesland & Bayern \\
Rechtsform & Körperschaft des öffentlichen Rechts / staatliche Einrichtung (bayerische Hochschule) \\
Sektor & öffentlich \\
Branche & Bildung und Forschung \\
Beschäftigte & über 1.000 \\
Personalvertretung & Personalrat \\
Gleichstellungsbeauftragte & ja \\
KI-Einsatz & KI-Tool zur Vorauswahl von Bewerbungen (Ranking) \\
Dienstleister außerhalb der EU & nein \\
Stelle & Wissenschaftliche Mitarbeiterin/wissenschaftlicher Mitarbeiter nach TV-L \\
Befristung & befristet auf 3 Jahre, Qualifizierungsstelle (Promotion) \\
\end{longtable}
\end{small}

\offen{Satz: Verweis auf Tabelle~\ref{tab:h1_prozess} (Prozess mit 6 Teilprozessen und 19 Tasks)}
\begin{small}
\begin{longtable}{@{}>{\raggedright\arraybackslash}p{4.2cm}>{\raggedright\arraybackslash}p{1.3cm}>{\raggedright\arraybackslash}p{9.0cm}@{}}
\caption[Prozess für Fall~1: Bewerbungsprozess wissenschaftliche Mitarbeiterstelle]{Prozess für Fall~1: Bewerbungsprozess wissenschaftliche Mitarbeiterstelle. Quelle: eigene Darstellung, Daten aus \repodatei{h1/input/prozess.json}.}\label{tab:h1_prozess}\\
\toprule \textbf{Teilprozess} & \textbf{Task} & \textbf{Bezeichnung} \\ \midrule \endfirsthead
\toprule \textbf{Teilprozess} & \textbf{Task} & \textbf{Bezeichnung} \\ \midrule \endhead
\bottomrule \endfoot
TP1 Stelle ausschreiben & T1.1 & Bedarf melden \\
 & T1.2 & Stellenprofil und Eingruppierung nach TV-L festlegen \\
 & T1.3 & Ausschreibung intern und extern veröffentlichen \\
TP2 Bewerbungen annehmen & T2.1 & Unterlagen über Bewerbungsportal empfangen \\
 & T2.2 & Eingang bestätigen \\
 & T2.3 & Bewerberdaten im System anlegen \\
TP3 KI-Vorauswahl & T3.1 & Unterlagen mit KI-Tool auswerten \\
 & T3.2 & Ranking erzeugen \\
 & T3.3 & Ergebnis durch Fachabteilung prüfen \\
 & T3.4 & Shortlist festlegen \\
TP4 Vorstellungsgespräche & T4.1 & Kandidaten einladen (inkl. Hinweis auf das Antragsrecht zur Beteiligung der Gleichstellungsbeauftragten) \\
 & T4.2 & Gespräche mit der Auswahlkommission führen (Gleichstellungsbeauftragte auf Antrag der Bewerbenden beteiligt) \\
 & T4.3 & Gespräche dokumentieren \\
TP5 Auswahl und Einstellung & T5.1 & Auswahlentscheidung treffen \\
 & T5.2 & Vertragsangebot erstellen \\
 & T5.3 & Arbeitsvertrag abschließen \\
TP6 Absagen und Löschen & T6.1 & Absagen versenden \\
 & T6.2 & Aufbewahrungsfrist abwarten \\
 & T6.3 & Bewerberdaten löschen \\
\end{longtable}
\end{small}

\offen{Satz: Auftrag an den Agenten (Prompt Stufe~1), Verweis auf \repodatei{h1/prompts/stufe1.md}. Entscheiden: Prompt vollständig abdrucken oder nur verlinken}

\subsubsection{Referenz}
\offen{Satz: Herleitung der Referenz aus der Referenzzuordnung (Anhang~\ref{anh:vorschriften}): es zählen Regelwerke mit mindestens einer Zeile der Art Zuordnung oder fehlender Schritt}
\begin{small}
\begin{longtable}{@{}>{\raggedright\arraybackslash}p{3.0cm}>{\raggedright\arraybackslash}p{3.0cm}>{\raggedleft\arraybackslash}p{3.2cm}>{\raggedleft\arraybackslash}p{3.8cm}@{}}
\caption[Referenz für H1 in Fall~1: 12 Regelwerke]{Referenz für H1 in Fall~1: 12 Regelwerke. Quelle: eigene Darstellung, Daten aus \repodatei{h1/input/soll_liste.csv}.}\label{tab:h1_referenz}\\
\toprule \textbf{Ebene} & \textbf{Regelwerk} & \textbf{Zeilen Zuordnung} & \textbf{Zeilen fehlender Schritt} \\ \midrule \endfirsthead
\toprule \textbf{Ebene} & \textbf{Regelwerk} & \textbf{Zeilen Zuordnung} & \textbf{Zeilen fehlender Schritt} \\ \midrule \endhead
\bottomrule \endfoot
EU & AI Act & 2 & 3 \\
EU & DSGVO & 6 & 1 \\
Bund & AGG & 5 & 0 \\
Bund & GG & 1 & 0 \\
Bund & NachwG & 1 & 0 \\
Bund & SGB IX & 1 & 3 \\
Bund & TzBfG & 1 & 0 \\
Bund & WissZeitVG & 1 & 0 \\
Land Bayern & BayDSG & 1 & 0 \\
Land Bayern & BayGlG & 2 & 0 \\
Land Bayern & BayPVG & 0 & 2 \\
Tarifvertrag & TV-L & 2 & 0 \\
\end{longtable}
\end{small}

\offen{Satz: neutrale Regelwerke (BayHIG) zählen weder als Treffer noch als Fehler; weitere vom Agenten genannte Regelwerke zählen als zusätzlich und senken die Genauigkeit}
\offen{Satz: Erfolgskriterium Trefferquote $\geq 0{,}9$, also mindestens 11 von 12 Regelwerken}
